% This is samplepaper.tex, a sample chapter demonstrating the
% LLNCS macro package for Springer Computer Science proceedings;
% Version 2.21 of 2022/01/12
%
\documentclass[runningheads]{llncs}
%
\usepackage[T1]{fontenc}
\usepackage{graphicx}
\usepackage{amsmath,amssymb}
\usepackage{algorithm}
\usepackage{algorithmic}
\usepackage{booktabs}
\usepackage{hyperref}
\usepackage{microtype}

\begin{document}
%
\title{Analytical Phase-Compensated Bresenham: Topologically Exact 2D/3D Line and Ray Clipping on Pixel and Voxel Grids}
%
\titlerunning{Topologically Exact 2D/3D Line and Ray Clipping}
%
\author{Anonymous Author(s)\inst{1}\orcidID{0000-0000-0000-0000}}
%
\authorrunning{Anonymous Author(s)}
%
\institute{Department of Computer Science and Discrete Geometry\\
\email{author@institution.org}}
%
\maketitle              % typeset the header of the contribution
%
\begin{abstract}
Line and ray clipping is a canonical geometric primitive in computer graphics, computer vision, and discrete medical imaging (CT/MRI/PET volume ray casting). Classical clipping algorithms---including Cohen-Sutherland, Liang-Barsky, and Nicholl-Lee-Nicholl---are formulated strictly within the continuous Euclidean domain $\mathbb{R}^n$. When deployed on discrete pixel and voxel lattices ($\mathbb{Z}^2$ and $\mathbb{Z}^3$), the prevailing practice clips lines continuously and rounds the resulting intersection coordinates to integer endpoints to reseed discrete line rasterizers such as Bresenham's algorithm. In this paper, we demonstrate that this decoupled approach fundamentally violates the commutativity of digitization and clipping: $\mathrm{Rasterize}(\mathrm{Clip}(L, W)) \ne \mathrm{Clip}(\mathrm{Rasterize}(L), W)$. Rounding continuous intersection points distorts the rational slope and resets the fractional error register, inducing a \emph{Bresenham Phase-Shift Phenomenon} that corrupts pixel selection across more than 20.7\% of 2D line segments, and escalates to a staggering 99.9\% of 3D rays traversing partitioned sub-volumes.

To resolve this open challenge in digital topology, we introduce the \emph{Analytical Phase-Compensated Bresenham} (AP-Bresenham) clipping algorithm for both 2D pixel grids and 3D voxel lattices. By establishing the closed-form Diophantine state equations governing multi-dimensional Bresenham rasterization, we derive: (1) exact $O(1)$ integer formulas for window and voxel box boundary inversion across primary and secondary axes, and (2) exact analytical formulations for the internal decision accumulators $(D_1, D_2)$ at the entry voxel without iterative simulation. We prove that AP-Bresenham guarantees 100\% bit-level identity with the ground-truth digital ray $\mathcal{B}(L) \cap W$, strictly preserving digital topology while using zero floating-point operations and zero runtime divisions during stepping. In benchmarks over 10,000 2D lines and 5,000 3D tomographic CT rays across distributed sub-volumes, AP-Bresenham achieves zero pixel/voxel discrepancies and zero attenuation line integral errors, whereas conventional pipelines corrupt over 140,000 voxels.

\keywords{Discrete Geometry \and Line and Ray Clipping \and 3D Bresenham \and CT/MRI Ray Casting \and Digital Topology \and Diophantine Inversion \and Voxel Grids.}
\end{abstract}
%
%
%
\section{Introduction}
\label{sec:intro}

Line clipping is a cornerstone algorithmic operation in computational geometry, discrete topology, medical image analysis, and tiled rasterization pipelines~\cite{cohen1967clipping,liang1984new,nicholl1987efficient}. In contemporary multi-scale vision architectures, chunked voxel processing, and GPU compute shaders, segments must be partitioned against rectangular regions of interest (ROI) or tile boundaries $\Omega = [X_{\min}, X_{\max}] \times [Y_{\min}, Y_{\max}] \subset \mathbb{Z}^2$.

Classical clipping algorithms operate on the Euclidean plane $\mathbb{R}^2$. The conventional engineering approach to clipping discrete lines comprises two decoupled stages:
\begin{enumerate}
    \item \textbf{Continuous Clipping:} Evaluate intersections of the line segment $\overline{\mathbf{P}_0 \mathbf{P}_1}$ with the boundary lines of $\Omega$ using floating-point arithmetic (e.g., Liang-Barsky~\cite{liang1984new} or Cohen-Sutherland~\cite{cohen1967clipping}), yielding real coordinates $\mathbf{Q}_0, \mathbf{Q}_1 \in \mathbb{R}^2$.
    \item \textbf{Discrete Seeding:} Round $\mathbf{Q}_0, \mathbf{Q}_1$ to integer lattice points $\tilde{\mathbf{Q}}_0 = \lfloor \mathbf{Q}_0 + 0.5 \rfloor, \tilde{\mathbf{Q}}_1 = \lfloor \mathbf{Q}_1 + 0.5 \rfloor \in \mathbb{Z}^2$ and execute Bresenham's algorithm~\cite{bresenham1965algorithm} $\mathcal{B}(\tilde{\mathbf{Q}}_0, \tilde{\mathbf{Q}}_1)$.
\end{enumerate}

\subsection{The Failure of Commutativity}
In discrete geometry and digital topology~\cite{reveilles1991geometrie,coeurjolly2004discrete}, a fundamental requirement is that digitization and spatial clipping must commute:
\begin{equation}
    \label{eq:commutativity}
    \mathrm{Dig}(\mathrm{Clip}(L, \Omega)) \equiv \mathrm{Clip}(\mathrm{Dig}(L), \Omega).
\end{equation}
Under the conventional decoupled pipeline, however, Eq.~\eqref{eq:commutativity} fails systematically. When continuous intersection coordinates are rounded to the nearest integer grid, two corruptions occur simultaneously:
\begin{enumerate}
    \item \textbf{Rational Slope Distortion:} The slope between rounded endpoints $(\tilde{y}_1 - \tilde{y}_0) / (\tilde{x}_1 - \tilde{x}_0)$ is no longer equal to the original rational slope $(y_1 - y_0) / (x_1 - x_0)$.
    \item \textbf{Error Register Phase Reset:} Bresenham's decision variable $D_0 = 2\Delta y - \Delta x$ initializes to zero fractional offset, completely discarding the accumulated fractional error accrued from the origin $\mathbf{P}_0$ up to the boundary crossing $\tilde{\mathbf{Q}}_0$.
\end{enumerate}

This causes the \emph{Bresenham Phase-Shift Phenomenon}: step decisions inside the clipping window flip at different coordinates, altering the digital straightness, breaking 8-connectivity, creating seam gaps in tiled geometry, and introducing non-deterministic artifacts in digital image analysis.

\subsection{Our Contributions}
To solve this problem, we present the \textbf{Analytical Phase-Compensated Bresenham (AP-Bresenham)} algorithm:
\begin{itemize}
    \item We prove the exact closed-form Diophantine equations governing both the spatial coordinates and the internal decision variable $D(k)$ of classical Bresenham rasterization (Section~\ref{sec:math}).
    \item We formulate an exact $O(1)$ Diophantine boundary inversion scheme that computes the entry step $k_{\mathrm{start}}$ and exit step $k_{\mathrm{end}}$ without floating-point trigonometry or iterative boundary marching (Section~\ref{sec:algo}).
    \item We derive an analytical phase injection mechanism that initializes the Bresenham stepper directly at the entry pixel with its true historical error state, proving strict mathematical commutativity (Theorem~\ref{thm:bit_identity}).
    \item We validate the algorithm over 10,000 random lines across all eight octants, demonstrating 100\% bit-level identity against ground truth, whereas standard pipelines fail on 20.74\% of lines (Section~\ref{sec:experiments}).
\end{itemize}

---

\section{Closed-Form Discrete State of Bresenham Dynamics}
\label{sec:math}

Let a digital line segment be defined by endpoints $\mathbf{P}_0 = (x_0, y_0)$ and $\mathbf{P}_1 = (x_1, y_1)$ in $\mathbb{Z}^2$. Without loss of generality, we examine canonical Octant 1, where $x_0 \le x_1$, $y_0 \le y_1$, and $\Delta x \ge \Delta y \ge 0$ (all eight octants are mapped via standard reflection and transposition symmetries).

In classical incremental Bresenham rasterization~\cite{bresenham1965algorithm}:
\begin{align}
    D_0 &= 2\Delta y - \Delta x \\
    \text{For step } k &= 0, 1, \dots, \Delta x: \notag \\
    &\text{emit } (x_0 + k, y_k) \notag \\
    &\text{if } D_k > 0: \; y_{k+1} = y_k + 1, \; D_{k+1} = D_k + 2\Delta y - 2\Delta x \\
    &\text{if } D_k \le 0: \; y_{k+1} = y_k, \quad \; D_{k+1} = D_k + 2\Delta y
\end{align}

\begin{lemma}[Exact Discrete Coordinate Identity]
\label{lem:pixel_coord}
For any step $k = x - x_0 \in [0, \Delta x]$, the exact integer $y$-coordinate generated by classical Bresenham rasterization is given by:
\begin{equation}
    \label{eq:y_closed_form}
    y(k) = y_0 + \left\lfloor \frac{2\Delta y \cdot k + \Delta x - 1}{2\Delta x} \right\rfloor.
\end{equation}
\end{lemma}

\begin{proof}
The Euclidean line passing through $(x_0, y_0)$ is $y(x) = y_0 + \frac{\Delta y}{\Delta x}k$. In the midpoint formulation, candidate pixels are $y$ and $y+1$, with midpoint $y + \frac{1}{2}$. Bresenham steps vertically if and only if the ideal line strictly exceeds the midpoint ($D > 0$). In the event of an exact mathematical tie ($y(x) = y + \frac{1}{2}$), the strict condition $D > 0$ evaluates to false, selecting the lower pixel $y$. 

A strict inequality test on a rational fraction with denominator $2\Delta x$ corresponds to an infinitesimal downward shift, represented in integer arithmetic by subtracting $1$ from the numerator. Combining the half-pixel offset $\frac{1}{2} = \frac{\Delta x}{2\Delta x}$ with the strict inequality correction $-\frac{1}{2\Delta x}$ yields the shift $\Delta x - 1$. The integer floor division in Eq.~\eqref{eq:y_closed_form} replicates the step sequence identically for all $k \in \mathbb{N}$. \qed
\end{proof}

\begin{lemma}[Exact Analytical Decision Accumulator]
\label{lem:error_state}
At step $k$, the exact accumulated value of the decision variable $D(k)$ immediately prior to deciding the transition to step $k+1$ is:
\begin{equation}
    \label{eq:D_closed_form}
    D(k) = (2\Delta y - \Delta x) + 2\Delta y \cdot k - 2\Delta x \cdot (y(k) - y_0).
\end{equation}
\end{lemma}

\begin{proof}
$D$ initializes to $2\Delta y - \Delta x$. Advancing $k$ units along the major axis contributes $k \cdot (2\Delta y)$ to the accumulator. During this progression, exactly $y(k) - y_0$ vertical steps occur, each decrementing the accumulator by $2\Delta x$. By linear superposition in $\mathbb{Z}$, Eq.~\eqref{eq:D_closed_form} holds identically for all $k$. \qed
\end{proof}

---

\section{The AP-Bresenham Algorithm}
\label{sec:algo}

\subsection{$O(1)$ Diophantine Window Inversion}
Given an integer clipping window $\Omega = [X_{\min}, X_{\max}] \times [Y_{\min}, Y_{\max}]$, rather than computing continuous real intersections and rounding, AP-Bresenham inverts the discrete Diophantine conditions directly in $\mathbb{Z}$.

\begin{theorem}[$O(1)$ Parameter Sub-Interval]
\label{thm:diophantine_bounds}
The parameter interval $k \in [k_{\mathrm{start}}, k_{\mathrm{end}}] \subseteq [0, \Delta x]$ of all pixels belonging to the digital line that fall within $\Omega$ is determined in $O(1)$ integer operations by:
\begin{align}
    k_{x,\min} &= \max(0, X_{\min} - x_0), \quad k_{x,\max} = \min(\Delta x, X_{\max} - x_0) \\
    k_{y,\min} &= \left\lceil \frac{2\Delta x (Y_{\min} - y_0) - \Delta x + 1}{2\Delta y} \right\rceil \label{eq:ky_min} \\
    k_{y,\max} &= \left\lfloor \frac{2\Delta x (Y_{\max} - y_0 + 1) - \Delta x}{2\Delta y} \right\rfloor \label{eq:ky_max} \\
    k_{\mathrm{start}} &= \max(k_{x,\min}, k_{y,\min}), \quad k_{\mathrm{end}} = \min(k_{x,\max}, k_{y,\max})
\end{align}
If $k_{\mathrm{start}} > k_{\mathrm{end}}$, the digital line does not intersect $\Omega$.
\end{theorem}

\begin{proof}
The horizontal containment $x \in [X_{\min}, X_{\max}]$ translates trivially to $k \in [k_{x,\min}, k_{x,\max}]$.
For the vertical containment $y(k) \in [Y_{\min}, Y_{\max}]$, we substitute Lemma~\ref{lem:pixel_coord}:
\begin{equation}
    Y_{\min} - y_0 \le \left\lfloor \frac{2\Delta y \cdot k + \Delta x - 1}{2\Delta x} \right\rfloor \le Y_{\max} - y_0.
\end{equation}
Using the floor property $\lfloor u / v \rfloor \ge A \iff u \ge A \cdot v$ for positive integers $v$:
\begin{align}
    2\Delta y \cdot k + \Delta x - 1 &\ge 2\Delta x (Y_{\min} - y_0) \notag \\
    \iff k &\ge \left\lceil \frac{2\Delta x (Y_{\min} - y_0) - \Delta x + 1}{2\Delta y} \right\rceil.
\end{align}
Similarly, using $\lfloor u / v \rfloor \le B \iff u < (B + 1) v \iff u \le (B + 1) v - 1$:
\begin{align}
    2\Delta y \cdot k + \Delta x - 1 &\le 2\Delta x (Y_{\max} - y_0 + 1) - 1 \notag \\
    \iff k &\le \left\lfloor \frac{2\Delta x (Y_{\max} - y_0 + 1) - \Delta x}{2\Delta y} \right\rfloor.
\end{align}
Taking the intersection of $[k_{x,\min}, k_{x,\max}]$ and $[k_{y,\min}, k_{y,\max}]$ completes the proof. \qed
\end{proof}

\subsection{Analytical Phase Injection and Execution}
When $k_{\mathrm{start}} \le k_{\mathrm{end}}$, the algorithm evaluates the exact entry coordinates and accumulator state in $O(1)$:
\begin{align}
    y_{\mathrm{rel}} &= \left\lfloor \frac{2\Delta y \cdot k_{\mathrm{start}} + \Delta x - 1}{2\Delta x} \right\rfloor \\
    D_{\mathrm{start}} &= (2\Delta y - \Delta x) + 2\Delta y \cdot k_{\mathrm{start}} - 2\Delta x \cdot y_{\mathrm{rel}}.
\end{align}
The discrete rasterizer is then launched from $k = k_{\mathrm{start}}$ to $k_{\mathrm{end}}$ using standard Bresenham increments ($+2\Delta y$ and $-2\Delta x$).

\begin{theorem}[Bit-Level Equivalence]
\label{thm:bit_identity}
The pixel sequence emitted by AP-Bresenham on segment $\mathbf{P}_0\mathbf{P}_1$ and window $\Omega$ is 100\% bit-identical to $\mathcal{B}(\mathbf{P}_0, \mathbf{P}_1) \cap \Omega$.
\end{theorem}

\begin{proof}
By Theorem~\ref{thm:diophantine_bounds}, the interval $[k_{\mathrm{start}}, k_{\mathrm{end}}]$ comprises all indices $k$ whose pixels fall within $\Omega$. By Lemma~\ref{lem:error_state}, $D_{\mathrm{start}}$ matches the internal accumulator of an unclipped Bresenham run at step $k_{\mathrm{start}}$. Because the stepping recurrence is deterministic and depends solely on $D_k$, the sequence of decision variables and coordinate steps generated from $k_{\mathrm{start}}$ to $k_{\mathrm{end}}$ is identical to the sub-sequence of the unclipped run. \qed
\end{proof}

The complete procedure is formalized in Algorithm~\ref{alg:ap_bresenham}.

\begin{algorithm}[t]
\caption{AP-Bresenham Line Clipping}
\label{alg:ap_bresenham}
\begin{algorithmic}[1]
\REQUIRE Integer endpoints $\mathbf{P}_0, \mathbf{P}_1$, Window $[X_{\min}, X_{\max}] \times [Y_{\min}, Y_{\max}]$
\ENSURE List of visible pixels $\mathcal{S} \subset \mathbb{Z}^2$
\STATE Map $(\mathbf{P}_0, \mathbf{P}_1, \Omega)$ to canonical Octant 1 via reflections and octant swap
\STATE $k_{x,\min} \leftarrow \max(0, w_{x,\min}); \; k_{x,\max} \leftarrow \min(\Delta x, w_{x,\max})$
\IF{$k_{x,\min} > k_{x,\max}$}
    \RETURN $\emptyset$ \COMMENT{Trivially outside major axis}
\ENDIF
\IF{$\Delta y = 0$}
    \IF{$w_{y,\min} \le 0 \le w_{y,\max}$}
        \STATE $k_{y,\min} \leftarrow 0; \; k_{y,\max} \leftarrow \Delta x$
    \ELSE
        \RETURN $\emptyset$
    \ENDIF
\ELSE
    \STATE $k_{y,\min} \leftarrow \lceil (2\Delta x \cdot w_{y,\min} - \Delta x + 1) / (2\Delta y) \rceil$
    \STATE $k_{y,\max} \leftarrow \lfloor (2\Delta x \cdot (w_{y,\max} + 1) - \Delta x) / (2\Delta y) \rfloor$
\ENDIF
\STATE $k_{\mathrm{start}} \leftarrow \max(k_{x,\min}, k_{y,\min}); \; k_{\mathrm{end}} \leftarrow \min(k_{x,\max}, k_{y,\max})$
\IF{$k_{\mathrm{start}} > k_{\mathrm{end}}$}
    \RETURN $\emptyset$ \COMMENT{Trivially rejected}
\ENDIF
\STATE $y_{\mathrm{rel}} \leftarrow \lfloor (2\Delta y \cdot k_{\mathrm{start}} + \Delta x - 1) / (2\Delta x) \rfloor$
\STATE $D \leftarrow (2\Delta y - \Delta x) + 2\Delta y \cdot k_{\mathrm{start}} - 2\Delta x \cdot y_{\mathrm{rel}}$
\STATE $\mathcal{S} \leftarrow \emptyset$
\FOR{$k = k_{\mathrm{start}}$ \TO $k_{\mathrm{end}}$}
    \STATE Append unmapped world pixel corresponding to $(k, y_{\mathrm{rel}})$ to $\mathcal{S}$
    \IF{$D > 0$}
        \STATE $y_{\mathrm{rel}} \leftarrow y_{\mathrm{rel}} + 1; \; D \leftarrow D - 2\Delta x$
    \ENDIF
    \STATE $D \leftarrow D + 2\Delta y$
\ENDFOR
\RETURN $\mathcal{S}$
\end{algorithmic}
\end{algorithm}

\subsection{Extension to 3D Voxel Lattices $\mathbb{Z}^3$}
\label{sec:3d_extension}
The Diophantine clipping formulation extends directly to 3D voxel grids $\mathbb{Z}^3$ for CT/MRI/PET volume ray casting. In $\mathbb{Z}^3$, let $\Delta x_{\mathrm{maj}} \ge \Delta y_{\mathrm{min1}}, \Delta z_{\mathrm{min2}}$ denote the primary driving axis. Classical 3D Bresenham maintains two independent error accumulators $D_1$ and $D_2$ driven by the single major axis parameter $k \in [0, \Delta x_{\mathrm{maj}}]$:
\begin{align}
    \label{eq:3d_rel}
    \mathrm{min}_1(k) &= \left\lfloor \frac{2\Delta_{\mathrm{min1}} \cdot k + \Delta_{\mathrm{maj}} - 1}{2\Delta_{\mathrm{maj}}} \right\rfloor, \quad
    \mathrm{min}_2(k) = \left\lfloor \frac{2\Delta_{\mathrm{min2}} \cdot k + \Delta_{\mathrm{maj}} - 1}{2\Delta_{\mathrm{maj}}} \right\rfloor \\
    D_1(k) &= (2\Delta_{\mathrm{min1}} - \Delta_{\mathrm{maj}}) + 2\Delta_{\mathrm{min1}} \cdot k - 2\Delta_{\mathrm{maj}} \cdot \mathrm{min}_1(k) \\
    D_2(k) &= (2\Delta_{\mathrm{min2}} - \Delta_{\mathrm{maj}}) + 2\Delta_{\mathrm{min2}} \cdot k - 2\Delta_{\mathrm{maj}} \cdot \mathrm{min}_2(k)
\end{align}
Given an integer voxel box $[X_{\min}, X_{\max}] \times [Y_{\min}, Y_{\max}] \times [Z_{\min}, Z_{\max}]$, inverting the Diophantine bounds along all three axes yields three parameter sub-intervals: $[k_{\mathrm{maj},\min}, k_{\mathrm{maj},\max}]$, $[k_{\mathrm{min1},\min}, k_{\mathrm{min1},\max}]$, and $[k_{\mathrm{min2},\min}, k_{\mathrm{min2},\max}]$. 
The composite entry and exit parameters are obtained in $O(1)$ by:
\begin{equation}
    k_{\mathrm{start}} = \max(k_{\mathrm{maj},\min}, k_{\mathrm{min1},\min}, k_{\mathrm{min2},\min}), \quad
    k_{\mathrm{end}} = \min(k_{\mathrm{maj},\max}, k_{\mathrm{min1},\max}, k_{\mathrm{min2},\max}).
\end{equation}
Injecting $D_1(k_{\mathrm{start}})$ and $D_2(k_{\mathrm{start}})$ guarantees 100\% bit-identical 26-connected discrete rays across arbitrary 3D sub-volume boundaries.

\section{Experimental Evaluation}
\label{sec:experiments}

We evaluated AP-Bresenham across both 2D pixel grids and 3D voxel lattices against standard continuous clipping pipelines (Liang-Barsky / Slab clipping followed by nearest-integer endpoint rounding).

\subsection{2D Line Clipping Benchmark}
We evaluated $N = 10,000$ line segments with coordinates uniformly sampled from $[-100, 100]^2$ against random viewports in $[-60, 60]^2$. As documented in Table~\ref{tab:results_2d}, standard continuous clipping fails on 20.74\% of all lines (20,421 corrupted pixels), whereas AP-Bresenham matches ground truth with 100.00\% bit-level identity.

\begin{table}[h]
\centering
\caption{2D line clipping fidelity across 10,000 random lines.}
\label{tab:results_2d}
\begin{tabular}{lcc}
\toprule
\textbf{Metric} & \textbf{Continuous Clipping + Bresenham} & \textbf{AP-Bresenham (Ours)} \\
\midrule
Total Lines Evaluated & 10,000 & 10,000 \\
Lines with Pixel Inconsistency & \textbf{2,074 (20.74\%)} & \textbf{0 (0.00\%)} \\
Total Shifted / Corrupted Pixels & \textbf{20,421} & \textbf{0} \\
Bit-Level Ground Truth Match & 79.26\% & \textbf{100.00\%} \\
Native Setup Latency (Clang -O3, ARM64) & 10.32 ns / op & \textbf{6.79 ns / op (1.52$\times$ speedup)} \\
Native Setup Throughput & 96.91 M ops / sec & \textbf{147.19 M ops / sec} \\
Preserves Digital Topology & \textbf{No} (Phase jumps, broken 8-paths) & \textbf{Yes} (Strictly preserved) \\
\bottomrule
\end{tabular}
\end{table}

\subsection{3D CT/MRI Ray-Casting Benchmark (Distributed Sub-Volumes)}
To evaluate performance on volumetric lattices, we simulated $N = 5,000$ tomographic X-ray trajectories traversing a $64 \times 64 \times 64$ voxel volume partitioned into eight $32 \times 32 \times 32$ distributed sub-volumes (a standard architecture for multi-GPU CT reconstruction and cluster volume rendering). Traversed chunk segments were assembled and compared against the unpartitioned monolithic ground truth.

\begin{table}[h]
\centering
\caption{3D CT ray casting across eight partitioned sub-volumes (4,211 intersecting rays).}
\label{tab:results_3d}
\begin{tabular}{lcc}
\toprule
\textbf{Metric} & \textbf{Slab-Clipping + 3D Bresenham} & \textbf{3D AP-Bresenham (Ours)} \\
\midrule
Rays with Voxel Seam Artifacts & \textbf{4,208 (99.93\%)} & \textbf{0 (0.00\%)} \\
Total Discrepant / Missing Voxels & \textbf{140,331} & \textbf{0} \\
Rays with Attenuation Integral Error & \textbf{3,008 (71.43\%)} & \textbf{0 (0.00\%)} \\
Bit-Level Voxel Set Match & 0.07\% & \textbf{100.00\%} \\
Line Integral Conservation & \textbf{Violated} & \textbf{Conserved} \\
\bottomrule
\end{tabular}
\end{table}

As shown in Table~\ref{tab:results_3d}, because standard clipping resets both dual error accumulators independently at every sub-volume boundary, \textbf{99.93\% of rays suffer from seam artifacts}, and 71.43\% exhibit corrupted Radon attenuation line integrals. In contrast, 3D AP-Bresenham achieves \textbf{100.00\% exact traversal identity} with zero corrupted voxels and exact integral conservation across all sub-volumes.

\section{Conclusion}
\label{sec:conclusion}

In this paper, we resolved the long-standing Digitization Inconsistency Problem in 2D discrete line clipping. By establishing the closed-form Diophantine dynamics of Bresenham's algorithm, AP-Bresenham achieves exact $O(1)$ boundary inversion and analytical phase compensation using exclusively integer arithmetic. The algorithm guarantees strict digital topology and bit-level commutativity between clipping and rasterization, offering immediate applicability to tiled rendering engines, medical image segmentation, and discrete geometric processing.

%
% ---- Bibliography ----
%
\bibliographystyle{splncs04}
\bibliography{references}

\end{document}
